% preprint/modules/lineage-correction.tex
% Shared-Horizon Preprint — Narrative Modules
% Module: Lineage Correction (Green Tile)

\section{Lineage Correction: Reversible Narrative Ecology}
\label{sec:lineage-correction}

A central contribution of the Shared-Horizon framework is the correction of the
green-tile lineage. Earlier interpretations classified the green tile as
\emph{symbolic ambiguity}, a category defined by unstable transitions, weak
expressive coherence, and non-reversible narrative behavior. This classification
was based on surface-level phenomenology rather than mathematical invariants.

Detailed analysis of expressive curvature, narrative reversibility, and
transition-map behavior demonstrates that the green tile is not symbolic
ambiguity. Instead, it corresponds to a distinct lineage with its own geometric
and narrative invariants: \emph{Reversible Narrative Ecology} (RNE).

\subsection{Evidence for Reversible Narrative Ecology}

The correction is supported by three mathematical observations:

\paragraph{1. Reversible Narrative Operators}

As shown in Section~\ref{sec:reversible-narrative-operators}, the green-tile
systems exhibit reversible narrative operators
\[
R : \mathcal{N} \rightarrow \mathcal{N}
\]
with smooth inverses and curvature-preserving behavior. Symbolic ambiguity lacks
reversibility and cannot support such operators.

\paragraph{2. Narrative Curvature Preservation}

Green-tile systems preserve narrative curvature:
\[
K_{\mathcal{N}}(R(p)) = K_{\mathcal{N}}(p).
\]

Symbolic ambiguity exhibits curvature drift, making curvature preservation
impossible. The presence of curvature-preserving transitions is a defining
feature of RNE.

\paragraph{3. Expressive–Narrative Compatibility}

Green-tile transitions commute with expressive transitions:
\[
R \circ T = T \circ R.
\]

This compatibility is mathematically incompatible with symbolic ambiguity, which
produces non-commutative expressive-narrative interactions. The compatibility
condition is a hallmark of reversible narrative systems.

\subsection{Collapse-Risk Analysis}

Collapse-risk geometry (Section~\ref{sec:collapse-risk-geometry}) further
supports the correction. Symbolic ambiguity exhibits high collapse curvature and
unstable collapse gradients, while green-tile systems remain within stable
collapse envelopes:
\[
C(R(p)) \le C(p).
\]

This stability is consistent with reversible narrative behavior and inconsistent
with symbolic ambiguity.

\subsection{Implications for the Mosaic}

Correcting the green tile has several implications for the Shared-Horizon
mosaic:

\begin{enumerate}
    \item It restores conceptual clarity by aligning the green tile with its
    true mathematical invariants.

    \item It ensures that narrative systems are compared using reversible
    operators, narrative geodesics, and curvature-preserving transitions.

    \item It prevents collapse between expressive and narrative manifolds by
    enforcing expressive–narrative compatibility.

    \item It strengthens the neutral membrane by ensuring that each tile
    corresponds to a mathematically coherent lineage.
\end{enumerate}

\subsection{Academic Context}

External academic references—including ANIMA Classic (Zenodo DOI:
10.5281/zenodo.22812349), the Anima GitHub repository, and NDH Meta-Systems
analysis—provide independent evidence for reversible narrative behavior. These
references are used strictly for scholarly comparison and do
